\documentclass[11pt,a4paper]{article}

\usepackage[margin=1in]{geometry}
\usepackage{booktabs}
\usepackage{multirow}
\usepackage{hyperref}
\usepackage{xcolor}
\usepackage{natbib}
\usepackage{array}
\usepackage{longtable}

\hypersetup{
  colorlinks=true,
  linkcolor=blue!60!black,
  citecolor=blue!60!black,
  urlcolor=blue!60!black
}

\title{%
  Thinking-Trace Analysis: \\
  Methodology for Surfacing Operational Improvements \\
  and Behavioral-Contract Candidates \\
  from Agent Planning Text%
}

\author{Sophie D.\ Neilson \\
  \small Independent Researcher \\
  \small ORCID: \href{https://orcid.org/0009-0001-2430-1743}{0009-0001-2430-1743}
}

\date{March 2026}

\begin{document}
\maketitle

\begin{abstract}
This paper records a methodology for analyzing AI agent thinking traces: the
extended planning text an agent produces before and during task execution. The
methodology uses a two-pass structure. The first pass is a piece-by-piece
temporal pass that extracts operational improvements, behavioral-contract
candidates, vocabulary coinages, and efficiency patterns per trace segment. The
second is a structural synthesis pass that ranks improvements by leverage and
attributes overhead to either the system (form and protocol gaps) or the agent
(judgment patterns). The central analytical question is whether observed
overhead is structural or behavioral, because the distinction determines the
appropriate remediation. The methodology is demonstrated on one trace from an
experimental inquest design task. Limitations include single-trace calibration,
analyst-dependent segmentation, and invisibility of absent reasoning.
\end{abstract}

\section{The Analytical Problem}

Agent thinking traces are a record of planning that disappears from the
accessible context once execution begins. The documents agents produce
(journals, inquests, accounts) record outputs: decisions made, artifacts
written, acceptance criteria evaluated. They do not record the planning
overhead that preceded those outputs. They do not show how many times a
direction was taken and abandoned, which constraints were discovered reactively
versus anticipated from the form, or where reasoning time was spent building
structure the form should have provided.

This overhead is recoverable only from thinking text. Reading a completed
artifact tells you what the agent decided. Reading the thinking trace tells you
the cost of that decision. It also tells you whether the cost was structural
(paid because a form or protocol lacked a gate) or behavioral (paid because the
agent's own approach had a pattern worth correcting).

That distinction is the central analytical question. It determines what kind of
fix is warranted. Structural overhead calls for form changes: a pre-design
gate, mandatory fields, an explicit re-entry protocol. Behavioral overhead
calls for agent advice: do X before Y, separate design from execution, read the
campaign document first. Conflating the two produces the wrong remediation.
Form changes get applied to things only the agent can catch. Agent advice gets
applied to things the form should simply require.

Thinking-trace analysis is the instrument for making this distinction reliably.

\subsection{Terminology}

The methodology was developed within a structured AI agent workflow system. The
following terms appear throughout the paper.

\paragraph{Agent roles.} Each role is an AI agent session with a defined scope
of responsibility.

\begin{table}[h]
\centering
\small
\caption{Agent roles referenced in this paper.}
\label{tab:roles}
\begin{tabular}{@{}lp{9.5cm}@{}}
\toprule
Role & Responsibility \\
\midrule
Researcher & Designs and executes experimental studies \\
Protocol Tester (PT) & Performs temporal, entry-by-entry analysis; identifies operational improvements and behavioral-contract candidates \\
Technowizard (TW) & Performs structural synthesis; ranks improvements by leverage and attributes overhead sources \\
Facilitating agent & Prepares and manages the study environment for each execution run (this role was coined during the analyzed trace) \\
Gardener & Manages dependency analysis and structural-change risk assessment \\
Synthesist & Performs evidence-weighting and cross-document pattern matching \\
\bottomrule
\end{tabular}
\end{table}

\paragraph{Document types.} Agents work within a structured document system.

\begin{table}[h]
\centering
\small
\caption{Document types referenced in this paper.}
\label{tab:doctypes}
\begin{tabular}{@{}lp{9.5cm}@{}}
\toprule
Type & Purpose \\
\midrule
Brief & A task specification with acceptance criteria that scopes a unit of work \\
Inquest & A structured investigation of an operational defect, following a defined sprint protocol \\
Journal & An analytical record organized as sequential entries, each documenting observations from a defined scope \\
Account & A post-execution record documenting what happened against the brief's acceptance criteria \\
Form & A structured template that defines required fields and checkpoints for a document type \\
TRACE & A per-session provenance record documenting what the agent did and when \\
Campaign & A coordination document that tracks bundled work items and priority constraints \\
\bottomrule
\end{tabular}
\end{table}

\paragraph{System concepts.}

\begin{itemize}
  \item \textbf{Form gate}: a required checkpoint in a form that must pass
    before the agent can proceed to the next step.
  \item \textbf{Taboo}: a behavioral contract that prevents locally rational
    but globally harmful actions. Identified by the X/Y/Z structure: X is the
    rational behavior, Y is the structural condition that permits it, Z is the
    globally bad outcome.
  \item \textbf{Thinking trace}: the extended planning text an AI agent
    produces before and during task execution. This text is not preserved in
    the agent's output documents. It is recoverable only from the model's
    reasoning output.
\end{itemize}

\section{The Two-Pass Structure and Why Each Pass Is Necessary}

\subsection{Pass 1: Piece-by-Piece (Protocol Tester)}

The first pass moves through the thinking trace in temporal order, treating
each reasoning segment as a unit of analysis. For each segment, it extracts:

\begin{itemize}
  \item \textbf{Operational improvements}: what the agent discovered reactively
    that a form gate should have surfaced proactively.
  \item \textbf{Taboo opportunities}: structural decision points where locally
    rational behavior could produce a globally bad outcome. These are
    candidates for behavioral contracts.
  \item \textbf{Vocabulary coinages}: terms the agent invents when the form has
    no prompt for the concept the agent is working with.
  \item \textbf{What worked well}: efficient behaviors the agent self-directed
    correctly, which the form does not require but should.
\end{itemize}

The piece-by-piece pass preserves temporal sequence. This matters because the
same constraint (a slug collision check, for example) may be identified in
entry~1 as a concern, again in entry~2 as a confirmed gap, and resolved in
entry~3, but at the wrong stage. The temporal record shows whether the agent
caught the constraint early or late, and whether the form prompted the check or
the agent generated it independently.

Vocabulary coinages are a diagnostic signal specific to this pass. When an
agent coins a term mid-execution (``vector-papering,'' ``facilitating agent,''
``annotation-resolution parity,'' ``gap-fire phase,'' ``designer's hypothesis
map''), the coinage marks a structural gap. The agent is inventing language
because the form has no field for the concept. Vocabulary that is not coined
but used fluently (from role files or prior forms) does not carry this signal.
Tracking coinages across the trace identifies where the form's conceptual
coverage ends.

\subsection{Pass 2: Structural Synthesis (Technowizard)}

The second pass reads the whole trace for system-level patterns. It does not
re-analyze individual segments. It builds on the piece-by-piece findings but
uses them as evidence for structural claims rather than as its primary output.

The structural pass extracts:

\begin{itemize}
  \item \textbf{Ranked system improvements}: what form or protocol changes
    would reduce overhead across all future executions of this type, ranked by
    leverage (how many sessions the fix improves and by how much).
  \item \textbf{Overhead source attribution}: for each improvement, whether the
    overhead source is the system (form lacks a gate) or the agent (judgment
    pattern worth correcting).
  \item \textbf{Specific agent advice}: per-behavior guidance concrete enough
    to change behavior on the next run. ``Do X before Y,'' not ``be more
    careful.''
  \item \textbf{Efficiency baselines}: what the agent did correctly that should
    not be disturbed. Examples: parallel reads at session start,
    format-calibration reads before writing, targeted verification before
    committing.
\end{itemize}

The structural pass also performs a correction function. Because it reads the
full trace after the piece-by-piece pass, it can identify where individual
segment analysis overclaimed. In the ouija-inquest-design trace, the
piece-by-piece pass raised a concern about H1/H3 near-collision (two
hypotheses both turning on timing). The structural pass partially corrected
this. The final hypothesis set (DEFAULT-UNSCOPED, TRIGGERS-MISSING-ENTRY,
PROTOCOL-INTENTIONAL-SPLIT) turns on different evidence types, and the overlap
concern applies more to an intermediate scenario design than to the final one.
Segment analysis compresses. The structural pass has context the segment did
not.

\subsection{Why Neither Pass Is Sufficient Alone}

The piece-by-piece pass produces a complete operational trace but no priority
ordering. A 30-item improvement list with no leverage ranking is not
actionable. The structural pass produces a ranked, attributed list but may miss
granular signal (vocabulary coinages, efficiency observations, individual
constraint-discovery moments) that does not survive compression into structural
claims.

The two passes have complementary coverage maps:

\begin{table}[h]
\centering
\small
\caption{Complementary coverage of the two analytical passes.}
\label{tab:coverage}
\begin{tabular}{@{}p{6.5cm}p{6.5cm}@{}}
\toprule
Piece-by-piece pass captures & Structural pass captures \\
\midrule
Temporal sequence of constraint discovery & Priority ordering by leverage \\
Vocabulary coinages as gap signals & Overhead source attribution \\
Per-segment ``what worked well'' & Specific agent advice \\
Taboo candidate identification & Leverage-ranked system improvements \\
Individual moment detail & Correction of segment overclaims \\
\bottomrule
\end{tabular}
\end{table}

A single-pass analysis using either method alone would produce findings, but
those findings would be either unprioritized (piece-by-piece alone) or
potentially disconnected from granular trace evidence (structural alone).

\section{Finding Taxonomy}

Four finding types emerged from applying this methodology to the
ouija-inquest-design trace. They are described here as a proposed taxonomy.
Whether these four types are sufficient or exhaustive across other traces and
task types is not yet known (see \S\ref{sec:single-trace}).

\subsection{Operational Improvements}

Specific form or process changes that would prevent a class of overhead.
Distinguished from agent advice by their addressability: an operational
improvement changes what a form requires; agent advice changes what an agent
does when the form is silent.

An operational improvement is correctly identified when the agent did the right
thing self-directed, without form prompting, and when the same behavior could
be made mechanical via a gate. The analysis guide structure (six components:
diagnostic reach, diagnostic assumption, paper-over, dead-end naming criterion,
gap-fire phase, baseline adjacency confirmation) is an operational improvement.
The Researcher invented the entire structure mid-execution. A future
facilitating agent reading only the form would have no prompt for any of it.
Making the analysis guide a mandatory form section with required fields
converts form-design work into field-population work. That is a straightforward
leverage gain.

\subsection{Taboo Opportunities}

Structural decision points where the agent is locally rational but where a
missing gate creates conditions for a globally bad outcome. Identified by the
X/Y/Z structure: X (locally rational behavior), Y (structural condition that
permits it), Z (globally bad outcome).

Taboo opportunities from thinking trace analysis are candidates, not registered
taboos. They route to a triage queue for the Protocol Tester, not directly to
the registry. The thinking trace is the evidence base for the candidate. The
piece-by-piece pass records where in the trace the candidate was identified.

From the ouija-inquest-design trace, three distinct taboo opportunities emerged
across the six entries:

\begin{itemize}
  \item \texttt{experimental-inquest-scenario-design}: covers the pre-design
    gate and analysis guide.
  \item \texttt{experimental-inquest-run-facilitation}: covers the facilitating
    agent role definition and run setup/teardown.
  \item \texttt{experimental-inquest-analysis-guide}: covers the analysis
    guide's mandatory fields specifically, if not bundled with scenario-design.
\end{itemize}

\subsection{Agent Behavior Signals}

Patterns in the agent's approach that no form change would address, because
they are judgment calls: sequencing decisions, awareness of when design and
execution are blurring, meta-habits about form submission timing.

These are identified when the agent's approach has a structural signature
across the trace. A one-time mistake is not a behavior signal. A recurrent
pattern is. The dissolution of the design/execution boundary (the agent
designing planted inquest content and simultaneously running its investigation
sprint in its head) occurs in Blocks~3, 4, and~5. The Campaign-read ordering
error (Campaign state read at close, not at open) occurs twice. A
three-occurrence pattern is a behavioral signal worth specific advice.

Agent behavior signals are addressable only via advice, not via form changes,
because they represent agent initiative applied in the wrong order, not a
missing form prompt.

\subsection{Efficiency Patterns}

What the agent did well that should not be disturbed. These are identified
explicitly because the default tendency of improvement-focused analysis is to
produce a list of deficits without recording the baseline. Without efficiency
patterns, a remediation pass has no way to distinguish ``this is fine, don't
touch it'' from ``this is also a problem.''

From the ouija-inquest-design trace, the efficiency patterns were: five
parallel reads at session open (correct bootstrapping), two
format-calibration reads before writing any artifact (produces
authentic-looking output), directory verification calls before committing
files, the ``neither artifact is wrong'' insight applied to the artifact
before archival (not just recorded in the account), and the analysis guide's
four-field annotation scheme (the highest-value output of the execution
despite being form-unguided).

\section{Findings from the Ouija-Inquest-Design Trace}

\subsection{The Dominant Overhead Pattern: Reactive Constraint Discovery}

The Researcher made the right call on every major constraint: artifact
existence, slug collision, immutability, path verification, baseline adjacency,
analysis guide structure. None of these were errors. Every one was discovered
reactively, as it arose in the execution. The EXPERIMENTAL\_INQUEST\_FORM
handed the Researcher a blank canvas when it should have handed them a
constraint checklist.

The cost is visible as reasoning cycles burned on scenarios that were always
going to fail. Multiple scenario designs were built and abandoned across the
trace. Two collapsed when the protocol files they referenced turned out not to
exist; others were abandoned after slug collision checks, internal consistency
failures, or newly discovered constraints. A single
\texttt{ls scaffolding/protocols/} check at the start could have killed the
first two in 30 seconds. All overhead before the final scenario (the INDEX.md
stale active entry) is attributable to the absence of a pre-design constraint
gate.

This is the highest-leverage system improvement: a mandatory pre-design
checklist (file existence, slug collision, immutability constraint, artifact
path declaration, hypothesis independence check) that converts reactive
discovery into a front-loaded gate.

\subsection{The Analysis Guide as the Most Important Missing Form Section}

The analysis guide's six components were invented by the Researcher in
reasoning and do not appear as required fields anywhere in the
EXPERIMENTAL\_INQUEST\_FORM. The six components are: diagnostic reach vs.\
standard hypothesis formation, diagnostic assumption vs.\ standard working
assumption, diagnostic paper-over vs.\ appropriate hedging, dead-end naming
criterion, gap-fire phase with calibrated round estimate, and baseline
adjacency confirmation.

This is the most consequential form gap in the trace. The analysis guide is the
primary output used by the facilitating agent during execution runs. A
facilitating agent reading only the form has no prompt to include any of these
components. The Researcher's work in this execution is the template. It should
become the form's required structure.

The gap-fire phase and self-correction path assessment are particularly
important omissions. Gap-fire phase specifies when in the investigation cycle
the diagnostic reach is most likely to fire. For the ouija trap, that is
Round~2 or later; early Round~1 reaches are less diagnostically useful.
Self-correction path assessment specifies whether the executing agent can
correct a diagnostic reach using available artifacts alone. For the ouija trap,
the answer is no: the required information is not in any artifact. Without
these two fields, an analysis guide is underspecified as a measurement tool.

\subsection{The Git Provenance Crack}

The planted inquest is dated 2026-03-22 but was committed 2026-03-23. The
analysis guide explicitly recommends \texttt{git log} as a standard
investigative tool for establishing timestamp sequences. An executing agent who
follows this recommendation will see that the planted inquest was created after
its stated date. This directly contradicts the TRACE-omission justification
(``predates the mandatory TRACE check'').

This is the single largest authenticity crack in the design. It was not flagged
anywhere in the artifacts and is not acknowledged in the analysis guide. The
form should require a git provenance assessment step: after artifacts are
designed, before they are written, state whether the planted artifact's commit
date will contradict its stated date, and if so, either note this in the
analysis guide or address it before execution.

\subsection{The Facilitating Agent Role}

The facilitating agent was invented entirely within the artifact. This role
creates the TRACE document before each run, copies the EXPERIMENTAL\_INQUEST to
the active directory, annotates observation fields after each run, and manages
planted artifact staging and teardown. The role is load-bearing. Without it, the
EXPERIMENTAL\_INQUEST does not run. It is currently defined only inside one
artifact. That makes it undiscoverable unless that artifact is read, and
unreliable unless every Researcher who designs an EXPERIMENTAL\_INQUEST invents
the same role in the same way.

The TRACE creation responsibility, the run-numbering convention
(\texttt{run[N]}), and the Campaign state timing (check at run start, not
post-artifact) all require a canonical protocol file, not an inference from one
execution's reasoning blocks.

\subsection{Vocabulary Coinages as Gap Signals}

Five terms were coined mid-execution where the form had no prompt for the
concept:

\begin{itemize}
  \item \textbf{Vector-papering}: treating reconstructed sequence as
    established intent. Papering over the gap between ``we can know the
    sequence'' and ``we can know what the completing agent intended.'' Named at
    the point the Researcher was articulating what the trap measures.
  \item \textbf{Facilitating agent}: the role that prepares the study
    environment before a run. Coined when the path conflict forced the
    Researcher to invent the role on the spot.
  \item \textbf{Annotation-resolution parity}: the baseline run resolved all
    annotation types; divergence in the experimental run on the same annotation
    types is the diagnostic signal. The sharpest statement of what the
    measurement tool is measuring.
  \item \textbf{Gap-fire phase}: the calibrated point in the investigation
    cycle where the diagnostic reach is most likely to occur. Named because the
    form had no field for it.
  \item \textbf{Designer's hypothesis map}: the intended hypothesis set,
    documented for the facilitating agent, with collapse-prone pairs flagged.
    Named when the H1/H3 near-collision concern surfaced and the analysis guide
    had no way to record it.
\end{itemize}

Each coinage marks the boundary of the form's conceptual coverage. A
vocabulary coinage is not a failure. It is evidence that the Researcher reached
something real that the form has not yet formalized.

\section{System Overhead vs.\ Agent Overhead}

\subsection{System-Generated Overhead}

System overhead arises when the agent performs the right behavior self-directed,
without form prompting. The fix is always a form or protocol change: adding the
required behavior as a gate. System overhead would be eliminated for all future
executions, regardless of the individual agent.

\begin{table}[h]
\centering
\small
\caption{System-generated overhead in the ouija-inquest-design trace.}
\label{tab:sysoverhead}
\begin{tabular}{@{}p{3.8cm}p{2.8cm}p{6cm}@{}}
\toprule
Overhead & Fix type & What changes \\
\midrule
Pre-design constraint discovery (two abandoned scenarios) & Pre-design gate & Form gains mandatory pre-design checklist \\[3pt]
Analysis guide structure invented from scratch & Mandatory form section & Analysis guide fields (six components) become required \\[3pt]
Artifact list grew from 3 to 7 items across execution & Artifact enumeration step & Complete artifact list locked before first file is written \\[3pt]
Facilitating agent role invented in artifact & Canonical protocol file & Protocol file for experimental-inquest facilitation \\[3pt]
Scenario pivot had no re-entry protocol & Pivot re-entry gate & If any gate fails, full design gate restarts \\[3pt]
Campaign state read at close, not open & Pre-execution checklist & Campaign read moved to first step of execution \\[3pt]
Verification step self-generated & Closure checklist & Form closure step: verify all deliverables exist at stated paths \\
\bottomrule
\end{tabular}
\end{table}

\subsection{Agent-Generated Overhead}

Agent overhead arises when the agent's own approach has a recurrent pattern
that the form cannot prevent, because it involves sequencing and judgment. The
fix is specific behavior advice. It changes what the agent does, not what the
form requires.

\paragraph{Design/execution boundary dissolution.} The Researcher was writing
planted artifact content and simultaneously running the investigation sprint in
its head, across three separate blocks. This produces artifacts whose content
is partly designed and partly discovered mid-write. That is a harder-to-track
boundary than a committed draft followed by a verification pass. The advice:
write a complete design draft, then run the verification sprint, then commit
the file. The sprint is a check, not a co-author.

\paragraph{Form update deferral.} The analysis guide structure (the
highest-value output of the execution) remains institutional knowledge in
reasoning blocks that are not committed anywhere load-bearing. The Researcher
identified this at close and did not open a brief for the form update before
archiving. The advice: when you invent structure that should be a form
requirement, open the brief before closing the session, not as a deferred
to-do.

\paragraph{``Neither artifact is wrong'' applied at close time.} The
Researcher identified at account-writing time that the strongest trap design
ensures no artifact looks obviously broken. Both are defensible individually;
the inconsistency only appears when the artifacts are read together against a
background question. This is a trap design heuristic, not a close-time
observation. The advice: apply it at design time, after locking the scenario.
Ask: ``is there an artifact in this set that looks obviously wrong?'' Redesign
until there is not.

\subsection{Why the Distinction Matters}

Applying agent advice to system-generated overhead produces brittle outcomes.
If a future agent has a different attention pattern, the reactive constraint
discovery returns. Applying form changes to agent-generated overhead is
category confusion. Forms cannot require judgment-sequencing.

The system/agent attribution is the analytical output that justifies the
two-pass structure. The piece-by-piece pass identifies the overhead. The
structural pass determines its source and the appropriate remediation type.

\section{Limitations}

\subsection{Absence Is Invisible}

Thinking-trace analysis can identify what the agent did, when the agent did it,
and whether the form prompted it. It cannot identify what the agent did not
think. If a constraint was never noticed (not discovered reactively, not
self-corrected, not mentioned in passing), the trace offers no signal. The git
provenance crack was flagged by both analysts because it was visible in the
artifact (the analysis guide recommends \texttt{git log}). An analogous crack
that produced no artifact trace would be invisible to both passes.

This is not a fixable limitation. It means that thinking-trace analysis is
reliable for the overhead it observes, and silent on the overhead it does not.

\subsection{Segmentation Is Analyst-Dependent}

The piece-by-piece pass requires the analyst to divide the thinking trace into
segments. The PT analysis used six blocks aligned to natural phase breaks:
setup through ``need to see archived structure''; artifact design through first
draft; artifact drafting through directory verification; artifact writing
through sprint execution; artifact writing through EXPERIMENTAL\_INQUEST
completion; account through close. A different analyst might segment more
finely (one segment per design decision) or more coarsely (design phase /
execution phase / close phase). The choice affects what the temporal sequence
shows. A finer segmentation surfaces more moment-level signals. A coarser
segmentation risks compressing away the reactive-discovery timing that matters
for system overhead attribution.

No specification currently exists for how to segment a thinking trace. This is
an open methodological question.

\subsection{Vocabulary Coinages Require Naming to Be Visible}

The five vocabulary coinages identified in \S\ref{sec:coinages} were visible
because the Researcher explicitly named them. Implicit vocabulary (concepts
that shaped behavior without being lexically marked) is not visible in the
trace. An agent who consistently treats artifact lists as complete when they
are partial, without ever naming the assumption, would not show a vocabulary
gap. It would show a recurring scope-expansion pattern that looks behavioral
rather than conceptual.

\subsection{``What Worked Well'' Is Structurally Underweighted}

The analytical attention in both passes naturally concentrates on gaps.
Efficiency patterns (\S3.4) are identified, but as a secondary category.
Over multiple trace analyses, this asymmetry could produce a picture of agent
behavior that is systematically negative: every trace reduced to a list of
what the agent got wrong. The efficiency baseline matters because it is the
thing to protect. It is the set of correct behaviors that form changes might
inadvertently discourage by adding overhead around them.

\subsection{Single-Trace Calibration}
\label{sec:single-trace}

The taxonomy and methodology in this study derive from one trace, one agent
(the Researcher), and one task type (experimental inquest design). The six
reasoning blocks in this trace map to five phases (setup, design, artifact,
execution, close). Whether that five-phase structure is intrinsic to agent
traces or specific to inquest design work is not determinable from one
instance. Whether the system/agent overhead distinction holds at similar
proportions for other role types or task types is unknown.

\section{Open Research}

\subsection{Overlap with Constructed-Environment Observation}

The ouija trap methodology (see the companion study~\citep{neilson2026canon})
and thinking-trace analysis are two instruments for observing agent behavior.
The ouija trap observes an executing agent navigating a structural gap during
genuine task execution. The agent has no meta-knowledge of the observation.
Thinking-trace analysis observes the planning and design agent explicitly
reasoning about what it is doing.

These are different behavioral surfaces. The ouija trap cannot see the
Researcher's design process. Thinking-trace analysis cannot see what the
executing agent does in the moment of encountering the gap during a sealed run.
They are likely to surface non-overlapping findings. The ouija trap reveals how
executing agents respond to a gap in their environment. Thinking-trace analysis
reveals how designing agents build gaps into the environment they create.

Whether they ever surface overlapping findings (the same behavioral pattern
visible from both directions) is an empirical question. A confirmed overlap
would be strong evidence for a stable behavioral pattern, not an artifact of
one observation method.

\subsection{Methodology Generalization to Other Roles}

The piece-by-piece + structural methodology was developed for a Researcher
trace. The Researcher's task (experimental inquest design) produces extended
reasoning about form constraints, artifact design, and scenario calibration. It
is a rich trace for operational improvement extraction.

Different roles would produce different trace structures. A Protocol Tester
trace would be rich in taboo candidate reasoning and behavioral contract
scoping. A Gardener trace would be rich in dependency analysis and
structural-change risk assessment. A Synthesist trace would be rich in
evidence-weighting and pattern-matching across source documents. Whether the
same four finding types (operational improvements, taboo opportunities, agent
behavior signals, efficiency patterns) appear across all role traces is
unknown. The two-pass structure (temporal extraction + structural synthesis)
should generalize, because the overhead source attribution question (system
vs.\ agent) applies to any role's reasoning.

\subsection{Vocabulary Coinage as a Systematic Form-Gap Detector}
\label{sec:coinages}

If vocabulary coinages reliably mark the boundaries of form conceptual
coverage, then tracking coinages across a corpus of thinking traces would
produce a map of where the form suite's coverage ends. This could be useful
for form maintenance: rather than waiting for a form gap to produce a visible
overhead pattern, coinage detection in traces could surface latent gaps before
they accumulate cost.

This depends on the reliability of the coinage signal. Some coinages are
productive. Annotation-resolution parity and designer's hypothesis map both
became required form fields. Some coinages mark invented infrastructure that
was needed for one execution and may not generalize (the specific facilitating
agent invocation pattern, for example). Distinguishing generalizable coinages
from one-off inventions may require multiple traces.

\subsection{Thinking-Trace Analysis as a Standing Practice}

This methodology was applied once to one trace. Its value scales with
frequency. The ouija-inquest-design trace covered six blocks and produced three
taboo candidates, five vocabulary coinages, seven system improvements
(\S5.1, consolidated from both passes), and three agent behavior signals, all
from a single task execution. A standing practice of trace analysis after
significant design executions would compound: form gaps identified in one
session become form requirements that reduce overhead in the next.

The question is resourcing. Trace analysis as practiced here required two full
analytical passes (PT and TW), each treating the trace at depth. This is not
trivially cheap. Whether a lighter-weight variant (a single-pass structural
read, or a PT pass only) would capture sufficient signal to justify the
reduction in rigor is worth testing.

\subsection{Source Documents and Data Availability}
\label{sec:data}

Source documents 1 through 4 are deposited alongside this paper as
supplementary files. They are also committed in the repository at
\texttt{papers/thinking-trace-analysis/sources/}.

\begin{table}[h]
\centering
\small
\caption{Source documents for this study.}
\label{tab:sources}
\begin{tabular}{@{}cp{4.5cm}p{7cm}@{}}
\toprule
\# & Filename & Contribution \\
\midrule
1 & {\small\raggedright\texttt{JOURNAL\_ouija-researcher-\allowbreak analysis\_\allowbreak 2026-03-23.md}} & Six-entry temporal pass: operational improvements, taboo candidates, vocabulary, and ``what worked well'' per segment \\[4pt]
2 & {\small\raggedright\texttt{JOURNAL\_technowizard-\allowbreak researcher-analysis\_\allowbreak 2026-03-23.md}} & Whole-trace structural read: six system improvements ranked by leverage, seven agent behavior advice items \\[4pt]
3 & {\small\raggedright\texttt{BRIEF\_ouija-inquest-\allowbreak design\_\allowbreak 2026-03-22.md}} & The task specification the Researcher executed. The thinking trace was the model's reasoning output during this execution; it is not stored in the brief \\[4pt]
4 & {\small\raggedright\texttt{STUDY\_STUB\_thinking-\allowbreak trace-analysis\_\allowbreak 2026-03-26.md}} & Scoping document that identified this as a study, not an account, and proposed the section structure \\[4pt]
5 & Ouija methodology study & Available from the author on request. Companion methodology study \\
\bottomrule
\end{tabular}
\end{table}

\bibliographystyle{plainnat}

\begin{thebibliography}{1}

\bibitem[Neilson(2026)]{neilson2026canon}
Neilson, S.~D. (2026).
\newblock Canon Suppression in Corpus-Loaded Assessment: A Two-Assessor Study.
\newblock Zenodo.
  \href{https://doi.org/10.5281/zenodo.22556876}{doi:10.5281/zenodo.22556876}.

\end{thebibliography}

\appendix

\section{Source Document Summaries}
\label{app:summaries}

Brief descriptions of each source document follow. The full texts of documents
1 through 4 are co-deposited as supplementary files.

\subsection{PT Analysis (Piece-by-Piece)}

The Protocol Tester journal (295 lines) walks through the Researcher's thinking
trace in six entries aligned to natural phase breaks. Each entry extracts four
categories: operational improvements, taboo opportunities, new vocabulary, and
what worked well. The journal closes with a whole-thing synthesis that rolls up
the per-entry findings into three taboo candidates and advice to the Researcher
on the git provenance crack, the H1/H3 near-collision, and the ``too easy''
calibration concern.

\subsection{TW Analysis (Structural)}

The Technowizard journal (430 lines) reads the same trace for system-level
patterns. It produces two ranked outputs: six system improvements ordered by
leverage (analysis guide structure, pre-design gate, artifact enumeration,
facilitating agent protocol, git provenance step, campaign read timing) and
seven items of specific agent advice (existence checks first, complete artifact
list, separate design from execution, read campaign first, submit form updates
before closing, apply the ``neither artifact is wrong'' heuristic at design
time, note the git provenance crack).

\subsection{Design Brief}

The design brief (155 lines) is the task specification for the
ouija-inquest-design execution. It defines the problem statement, acceptance
criteria, and context the Researcher worked from. The thinking trace itself was
the model's reasoning output during this execution. It is not stored in the
brief file. The PT and TW journals (source documents 1 and 2) quote the trace
extensively across their analyses. The trace spans six reasoning blocks
covering scenario design, artifact drafting, sprint execution, and task
closure.

\subsection{Study Stub}

The study stub (62 lines) is the scoping document that identified the two
analytical journals as constituting a methodology rather than a record. It
proposed the section structure used in this paper and named the research
question: what does extended thinking text reveal about agent planning and
design behavior, and can systematic multi-pass analysis reliably surface
operational improvements and structural taboo candidates?

\end{document}
